\documentclass[11pt,a4paper]{article}

\usepackage[utf8]{inputenc}
\usepackage[T1]{fontenc}
\usepackage[margin=1in]{geometry}
\usepackage{graphicx}
\usepackage{hyperref}
\usepackage{listings}
\usepackage{booktabs}
\usepackage{xcolor}
\usepackage{enumitem}
\usepackage{fancyhdr}

\hypersetup{
  colorlinks=true,
  linkcolor=blue!60!black,
  urlcolor=blue!60!black,
  citecolor=blue!60!black
}

\definecolor{codebg}{RGB}{245,245,245}
\definecolor{codeframe}{RGB}{220,220,220}

\lstdefinestyle{projectstyle}{
  backgroundcolor=\color{codebg},
  frame=single,
  rulecolor=\color{codeframe},
  basicstyle=\ttfamily\small,
  breaklines=true,
  columns=fullflexible,
  keepspaces=true,
  showstringspaces=false,
  tabsize=2
}

\lstset{style=projectstyle}

\pagestyle{fancy}
\fancyhf{}
\rhead{AIAM Technology LLC}
\lhead{Project Documentation}
\cfoot{\thepage}

\title{\textbf{AIAM Technology LLC}\\[0.4em]
\large Website Project Documentation}
\author{AIAM Technology LLC}
\date{\today}

\begin{document}

\maketitle
\thispagestyle{empty}
\vfill
\begin{center}
\textit{Formal website preview --- full launch planned for June 2026}
\end{center}
\newpage

\tableofcontents
\newpage

%==============================================================================
\section{Introduction}
%==============================================================================

This document describes the AIAM Technology LLC marketing website: its purpose,
architecture, source layout, runtime behavior, deployment on \textbf{Netlify},
and contact-form handling through \textbf{Formspree}.

The site is a static, client-side web application built with HTML, CSS, and
JavaScript. It presents four primary sections (Home, About, Technology, and
Contact) within a single page, using hash-based view switching so navigation
feels like separate tabs without full page reloads.

\subsection{Project Goals}

\begin{itemize}[noitemsep]
  \item Present AIAM's machine-side automation capabilities to manufacturers.
  \item Provide a professional landing experience aligned with the approved
        visual mockup (\texttt{src/img/aiam007.png}).
  \item Offer a working contact channel before the formal website launch.
  \item Deploy easily as static files on Netlify with Git-based continuous
        deployment.
\end{itemize}

\subsection{Repository}

\begin{itemize}[noitemsep]
  \item \textbf{GitHub:} \url{https://github.com/ASilvaGSW/AIAM.git}
  \item \textbf{Default branch:} \texttt{main}
  \item \textbf{Contact email (public):} \texttt{kenjin@aiam-tech.com}
  \item \textbf{Operating region:} Mission, Texas / Reynosa, Mexico
\end{itemize}

%==============================================================================
\section{Technology Stack}
%==============================================================================

\begin{table}[h]
  \centering
  \begin{tabular}{@{}ll@{}}
    \toprule
    \textbf{Layer} & \textbf{Technology} \\
    \midrule
    Markup       & HTML5 (semantic sections, ARIA labels) \\
    Styling      & CSS3 (Flexbox, Grid, custom properties, media queries) \\
    Scripting    & Vanilla JavaScript (ES5-compatible IIFE) \\
    Typography   & Google Fonts --- Inter \\
    Hosting      & Netlify (static site hosting) \\
    Form backend & Formspree (\url{https://formspree.io}) \\
    Version control & Git / GitHub \\
    \bottomrule
  \end{tabular}
  \caption{Technology stack summary}
\end{table}

No build step, package manager, or server-side runtime is required. The site
can be opened locally by serving the repository root or by opening
\texttt{index.html} through a local HTTP server (recommended for Formspree
testing).

%==============================================================================
\section{Project Structure}
%==============================================================================

The publish root for Netlify is the repository root. \texttt{index.html} must
remain at the top level so Netlify serves it as the site entry point.

\begin{lstlisting}
AIAM/
|-- index.html          # Site entry point and all page views
|-- css/
|   `-- styles.css      # Global layout, components, responsive rules
|-- js/
|   `-- main.js         # Navigation, view switching, Formspree submit
|-- src/
|   `-- img/            # Brand and content images
|       |-- Logo_5_19_2026.png
|       |-- aiam007.png           # Original design mockup (reference)
|       |-- hero_machine.png      # Hero section image
|       |-- card_wire.png         # Wire Harness card thumbnail
|       |-- card_vision.png       # AI Vision card thumbnail
|       `-- card_backup.png       # Power Backup card thumbnail
`-- docs/
    `-- project-documentation.tex # This document
\end{lstlisting}

\subsection{File Responsibilities}

\begin{description}[style=nextline]
  \item[\texttt{index.html}]
    Defines the header, four view sections (\texttt{\#view-home},
    \texttt{\#view-about}, \texttt{\#view-technology}, \texttt{\#view-contact}),
    the shared footer value strip, and the contact form markup.

  \item[\texttt{css/styles.css}]
    Implements the visual design, sticky header, hero diagonal image clip,
    feature cards, inner-page layouts, mobile navigation, footer stickiness,
    and the fit-to-viewport layout for large desktop screens.

  \item[\texttt{js/main.js}]
    Handles mobile menu toggle, client-side view routing via \texttt{data-view}
    attributes, URL hash updates, and asynchronous Formspree form submission.

  \item[\texttt{src/img/}]
    Static image assets. Hero and card images were derived from the approved
    mockup for visual consistency.
\end{description}

%==============================================================================
\section{Site Architecture and Navigation}
%==============================================================================

\subsection{Single-Page, Multi-View Model}

Instead of separate HTML files per section, the application uses a
\textbf{view-switching} pattern:

\begin{enumerate}[noitemsep]
  \item Each major section is a \texttt{<section class="view">} element with a
        \texttt{data-view} attribute (\texttt{home}, \texttt{about},
        \texttt{technology}, \texttt{contact}).
  \item Only the view with class \texttt{active} is displayed.
  \item Navigation links, the logo, and in-page CTAs carry matching
        \texttt{data-view} attributes.
  \item JavaScript intercepts clicks, prevents default navigation, and toggles
        the active view.
  \item The browser URL hash is updated (e.g.\ \texttt{\#contact}) using
        \texttt{history.replaceState}, enabling direct links and refresh
        persistence.
\end{enumerate}

\subsection{Page Content Overview}

\begin{table}[h]
  \centering
  \small
  \begin{tabular}{@{}p{2.2cm}p{10cm}@{}}
    \toprule
    \textbf{View} & \textbf{Content} \\
    \midrule
    Home &
    Hero with company headline, description, ``Formal Website Coming June 2026''
    badge, tagline, contact row, hero machine image, and three feature cards
    (Wire Harness, AI Vision, Machine Power Backup). \\
    About &
    Company mission, regional presence, tagline, and four value propositions
    with icons. \\
    Technology &
    Expanded descriptions of the three core technology areas with larger cards. \\
    Contact &
    Contact information and a message form (name, email, message) integrated
    with Formspree. \\
    \bottomrule
  \end{tabular}
  \caption{View sections}
\end{table}

\subsection{Shared Footer}

A dark value strip appears on every view, listing:

\begin{itemize}[noitemsep]
  \item Engineered for Reliability
  \item Modular by Design
  \item Built for Manufacturability
  \item Practical Support, Real Results
\end{itemize}

The layout uses a flex column on \texttt{body} so the footer remains at the
bottom of the viewport on desktop, even when inner-page content is shorter
than the screen height.

%==============================================================================
\section{Responsive Design and Layout Behavior}
%==============================================================================

\subsection{Breakpoints}

\begin{itemize}[noitemsep]
  \item \textbf{Desktop / laptop ($\geq$ 961px):} Two-column hero, three-column
        card grids, horizontal navigation.
  \item \textbf{Tablet ($\leq$ 960px):} Stacked hero, single-column cards and
        inner-page grids.
  \item \textbf{Mobile ($\leq$ 720px):} Hamburger menu, increased side padding,
        vertical card layout, full-width submit button.
  \item \textbf{Small mobile ($\leq$ 400px):} Further reduced typography and
        padding.
\end{itemize}

\subsection{Fit-to-Viewport Mode (Large Screens)}

On viewports with \texttt{min-width: 961px} and \texttt{min-height: 740px},
additional CSS rules apply:

\begin{itemize}[noitemsep]
  \item \texttt{body} is locked to \texttt{100vh} with \texttt{overflow: hidden}.
  \item The Home view compresses vertical spacing so hero, cards, header, and
        footer fit without scrolling.
  \item Inner views may scroll internally if content exceeds available height.
\end{itemize}

On smaller or shorter screens, normal document scrolling is used instead.

%==============================================================================
\section{JavaScript Behavior (\texttt{main.js})}
%==============================================================================

\subsection{Mobile Navigation}

The hamburger button toggles the \texttt{open} class on \texttt{\#mainNav} and
updates \texttt{aria-expanded} for accessibility.

\subsection{View Routing}

The \texttt{showView(name)} function:

\begin{itemize}[noitemsep]
  \item Activates the matching \texttt{.view} section.
  \item Updates active state on \texttt{.nav-link} elements.
  \item Resets scroll on the newly shown view.
  \item Closes the mobile menu.
  \item Writes the hash fragment to the URL.
\end{itemize}

On initial load, if the URL contains a valid hash (e.g.\ \texttt{\#technology}),
that view is opened automatically.

\subsection{Contact Form Submission}

The contact form uses the Fetch API to POST to Formspree without leaving the
page. See Section~\ref{sec:formspree} for backend details.

%==============================================================================
\section{Hosting on Netlify}
%==============================================================================

Netlify is a static hosting platform well suited to this project because the
site consists entirely of HTML, CSS, JavaScript, and image assets with no
server-side code.

\subsection{Why \texttt{index.html} Is at the Repository Root}

Netlify (and most static hosts) treat the configured \textbf{publish directory}
as the web root. The file \texttt{index.html} in that directory becomes the
default document for \texttt{/}. Placing \texttt{index.html} at the repository
root ensures:

\begin{itemize}[noitemsep]
  \item Correct resolution of relative paths such as \texttt{css/styles.css},
        \texttt{js/main.js}, and \texttt{src/img/...}.
  \item No extra redirect or rewrite rules are needed for the home page.
  \item Asset URLs remain simple and portable.
\end{itemize}

\subsection{Recommended Netlify Configuration}

\begin{table}[h]
  \centering
  \begin{tabular}{@{}ll@{}}
    \toprule
    \textbf{Setting} & \textbf{Value} \\
    \midrule
    Publish directory & \texttt{/} (repository root) \\
    Build command     & \emph{(none)} \\
    Production branch & \texttt{main} \\
    \bottomrule
  \end{tabular}
  \caption{Typical Netlify site settings for this project}
\end{table}

Because there is no build pipeline, Netlify deploys files as-is after each
push to the connected Git branch.

\subsection{Deployment Workflow (Git Integration)}

\begin{enumerate}
  \item Connect the Netlify site to the GitHub repository
        \url{https://github.com/ASilvaGSW/AIAM.git}.
  \item Set the production branch to \texttt{main}.
  \item Confirm the publish directory is the repository root.
  \item Push commits to \texttt{main}; Netlify triggers an automatic deploy.
  \item Review the deploy log in the Netlify dashboard if a deployment fails.
  \item Optionally assign a custom domain (e.g.\ \texttt{aiam-tech.com}) under
        \textbf{Domain management} and configure DNS with Netlify or an external
        registrar.
\end{enumerate}

\subsection{Optional: \texttt{netlify.toml}}

Although not required, a minimal \texttt{netlify.toml} at the repository root
can document deployment settings explicitly:

\begin{lstlisting}
[build]
  publish = "."
  command = ""

[[headers]]
  for = "/*"
  [headers.values]
    X-Frame-Options = "DENY"
    X-Content-Type-Options = "nosniff"
    Referrer-Policy = "strict-origin-when-cross-origin"
\end{lstlisting}

This file is optional; equivalent settings can be configured entirely in the
Netlify UI.

\subsection{Local Preview Before Deploy}

To approximate the Netlify environment locally:

\begin{lstlisting}[language=bash]
# Using Python (if installed)
python -m http.server 8080

# Or using Node npx
npx serve .
\end{lstlisting}

Then open \url{http://localhost:8080}. Use HTTP (not \texttt{file://}) when
testing Formspree submissions.

\subsection{Netlify Considerations for This Project}

\begin{itemize}
  \item \textbf{HTTPS:} Netlify provides free TLS certificates; Formspree works
        over HTTPS in production.
  \item \textbf{SPA-style hashes:} Hash routes (\texttt{\#about}, etc.) do not
        require special Netlify redirects because routing is handled client-side.
  \item \textbf{Form submissions:} Forms are sent directly from the browser to
        Formspree; Netlify does not process form data unless Netlify Forms is
        enabled separately (not used in this project).
  \item \textbf{Caching:} After CSS/JS updates, users may need a hard refresh
        (\Ctrl+Shift+R) if assets appear stale due to browser cache.
\end{itemize}

%==============================================================================
\section{Form Collection with Formspree}
\label{sec:formspree}
%==============================================================================

Formspree is a third-party service that receives HTML form submissions and
forwards them by email without requiring a custom backend.

\subsection{Endpoint}

The contact form posts to:

\begin{center}
\url{https://formspree.io/f/xvzyywkz}
\end{center}

This endpoint is configured in \texttt{index.html}:

\begin{lstlisting}
<form class="contact-form" id="contactForm"
      action="https://formspree.io/f/xvzyywkz"
      method="POST" novalidate>
\end{lstlisting}

\subsection{Form Fields}

\begin{table}[h]
  \centering
  \begin{tabular}{@{}lll@{}}
    \toprule
    \textbf{Field name} & \textbf{Input type} & \textbf{Required} \\
    \midrule
    \texttt{name}    & text     & Yes \\
    \texttt{email}   & email    & Yes \\
    \texttt{message} & textarea & Yes \\
    \bottomrule
  \end{tabular}
  \caption{Contact form fields sent to Formspree}
\end{table}

Field names must match Formspree expectations. The \texttt{email} field is
used by Formspree for reply-to addressing and validation.

\subsection{Submission Method (AJAX via Fetch)}

Rather than a traditional full-page POST, \texttt{main.js} intercepts submit
events and sends data asynchronously:

\begin{lstlisting}
fetch(form.action, {
  method: "POST",
  body: new FormData(form),
  headers: { Accept: "application/json" }
})
\end{lstlisting}

Benefits of this approach:

\begin{itemize}[noitemsep]
  \item The user stays on the Contact view (consistent with SPA navigation).
  \item Success and error messages appear inline below the submit button.
  \item The submit button shows a ``Sending\ldots'' state and is disabled
        during the request.
\end{itemize}

\subsection{User Feedback}

\begin{itemize}[noitemsep]
  \item \textbf{Success:} ``Thanks --- your message has been sent. We'll get
        back to you soon.'' The form is reset.
  \item \textbf{Validation error:} Browser-native validation via
        \texttt{checkValidity()} before submit.
  \item \textbf{Formspree error:} Server error messages are displayed when
        Formspree returns JSON error details.
  \item \textbf{Network error:} User is prompted to retry or email directly.
\end{itemize}

\subsection{Formspree Account Setup}

\begin{enumerate}
  \item Create or sign in to a Formspree account at \url{https://formspree.io}.
  \item Create a form and note the form ID (\texttt{xvzyywkz} in this project).
  \item Set the notification email address where submissions should be delivered.
  \item On the \textbf{first real submission}, Formspree typically sends a
        confirmation email to verify the form endpoint --- approve it to activate
        delivery.
  \item Optionally enable spam filtering (reCAPTCHA, honeypot) in the Formspree
        dashboard for production hardening.
\end{enumerate}

\subsection{Testing Checklist}

\begin{itemize}
  \item Serve the site over HTTP/HTTPS (not \texttt{file://}).
  \item Submit a test message from the Contact view.
  \item Confirm receipt in the Formspree dashboard and notification inbox.
  \item Test invalid email and empty fields to verify client-side validation.
  \item Test from the live Netlify URL after deployment.
\end{itemize}

\subsection{Security and Privacy Notes}

\begin{itemize}[noitemsep]
  \item The Formspree endpoint URL is public in client-side code; this is normal
        for static sites. Abuse mitigation should be configured in Formspree.
  \item Do not embed API secrets in frontend code; only the public form endpoint
        is used.
  \item Consider adding a privacy notice if required by jurisdiction or company
        policy before the formal launch.
\end{itemize}

%==============================================================================
\section{Maintenance and Future Work}
%==============================================================================

\subsection{Routine Updates}

\begin{itemize}[noitemsep]
  \item Edit content in \texttt{index.html} for copy changes.
  \item Edit \texttt{css/styles.css} for layout and visual updates.
  \item Replace images in \texttt{src/img/} keeping filenames or updating HTML
        references accordingly.
  \item Commit and push to \texttt{main} to trigger Netlify redeployment.
\end{itemize}

\subsection{Planned for Formal Launch (June 2026)}

\begin{itemize}[noitemsep]
  \item Expanded content and dedicated pages if moving beyond single-page views.
  \item Custom domain and production SEO metadata.
  \item Analytics integration (e.g.\ Netlify Analytics, Plausible, or GA4).
  \item Enhanced Formspree configuration (spam protection, autoresponders).
  \item Removal of ``Formal Website Coming June 2026'' preview messaging.
\end{itemize}

\subsection{Troubleshooting}

\begin{description}[style=nextline]
  \item[Styles or scripts not loading on Netlify]
    Verify \texttt{index.html} is at the publish root and paths use
    \texttt{css/...} and \texttt{js/...} without \texttt{../}.

  \item[Form submissions fail locally]
    Use a local HTTP server; Formspree may reject requests from
    \texttt{file://} origins.

  \item[Form works locally but not on Netlify]
    Confirm the Formspree form is activated and the site URL is allowed if
    domain restrictions are enabled in Formspree.

  \item[Footer not at bottom on desktop]
    Ensure \texttt{body} flex layout rules in \texttt{styles.css} are intact
    and not overridden by custom edits.

  \item[Mobile layout too tight]
    Check \texttt{@media (max-width: 720px)} rules; container horizontal
    padding is controlled by \texttt{.container}.
\end{description}

%==============================================================================
\section{Conclusion}
%==============================================================================

The AIAM Technology LLC website is a lightweight static application optimized
for Netlify deployment and Formspree-powered contact collection. Its structure
keeps operational complexity low while supporting a polished multi-section
experience on desktop and mobile. This document should serve as the primary
reference for developers maintaining, deploying, and extending the project
through the June 2026 formal launch.

\end{document}
